%% chapters/19-spiked-covariance.tex
\chapter{Spiked Covariance Models}
\label{ch:spiked}

\whereWeAre{%
The Marchenko--Pastur picture describes the eigenvalues of a sample covariance matrix when the population is isotropic ($\Sigma = I_p$). But what if the population has \emph{structure}, a low-rank ``signal'' on top of isotropic noise? That structure either does or does not produce visible eigenvalue spikes above the bulk; this chapter sets up the question and motivates the answer.%
}

\PLACEHOLDER{Mostly new chapter. Reference: Yale Stat 615 Lecture 9, Vershynin Ch.~7--8.}

\section{The spiked model: noise plus low-rank signal}
\PLACEHOLDER{$\Sigma = I_p + \sum_k \theta_k v_k v_k^T$. PCA motivation. Why this is the canonical model in modern data science.}

\section{Examples in the wild}
\PLACEHOLDER{Factor models in finance, latent semantic analysis, gene expression PCA, deep learning Hessian.}

\section{What a spike "looks like" in simulations}
\PLACEHOLDER{Histograms with spike clearly above the MP bulk; scatter of (sample top eigenvalue, true spike strength).}

\section{Brief warm-ups}
\PLACEHOLDER{3--5 short questions.}

\chapterRecap{%
\begin{itemize}
  \item \PLACEHOLDER{Spiked model = isotropic noise + structured low-rank signal.}
  \item \PLACEHOLDER{Spikes show up in the spectrum as outliers above the MP bulk.}
  \item \PLACEHOLDER{Whether they show up depends on signal strength, BBP transition in Chapter 20.}
\end{itemize}%
}

\section*{Exercises}
\PLACEHOLDER{8--15 longer exercises.}
